\SourceLeaf{161}$m, m', m'', m'''$ etc. selbst, doch ihr Verhältniss (i. e. das Verhältniss der Zahlen,
wodurch sie ausgedrückt werden, indem im Fall der Heterogeneität beliebig Einheiten für
jedes zum Grunde gelegt werden) kennen. Mittel zu einiger Schätzung wenigstens wird man bei
Mutterwitz fast immer finden; wo aber alle Kenntniss gänzlich fehlt, da liegt es doch
wahrhaftig nicht an der Methode selbst, dass sie ihre Dienste versagen muss. \textls{Sachkenntniss}
kann bei Anwendung der M. d. kl. Q. niemals erlassen werden. So wird es also eben von der
\textls{Sachkenntniss} jedesmal abhängen, in wie fern man bei \textls{Linearmessungen} den zu befürchtenden
mittlern Fehler für kleinere und grössere Distanzen gleich oder ungleich, und im letztern
Falle, wie man ihn von der Grösse abhängig setzen soll. Exempli gratia bei kleinen
Distanzen, wo z. B. mit einer 3 Meter langen Messruthe gemessen ist, wird, wenn die Distanz
selbst nur \textls{sehr wenige} Meter beträgt, der Fehler ganz oder fast allein von der Theilung
abhängen; muss aber der Massstab viele hundert male umgeschlagen werden, so tritt eine ganz
verschiedene Fehlerquelle ins Spiel.

\begin{center}
. . . . . . . . . . . .
\end{center}
